\documentclass[10pt,a4paper]{amsart}
\usepackage[margin=19mm]{geometry}
\usepackage{amsmath,amssymb,mathtools}
\usepackage[scr=rsfs,cal=euler]{mathalfa}
\usepackage{xfrac}
\usepackage[unicode,hidelinks,bookmarks=false]{hyperref}
\newcommand{\ie}{i.e.\ }
\newcommand{\cf}{cf.\ }
\newcommand{\resp}{respectively\ }
\newcommand{\Subsection}{\subsection}
\providecommand{\nobreakdash}{\nobreak\hbox{-}}
\setlength{\emergencystretch}{2em}
\multlinegap0pt
\usepackage{bm}
\usepackage{bbm}
\usepackage{dsfont}
\usepackage{xspace}
\usepackage{cancel}
\usepackage{mathtools}
\usepackage{amsthm}
\newtheorem{lem}{Lemma}
\newtheorem{prop}[lem]{Proposition}
\newcommand{\e}{e}
\renewcommand{\ge}{\geqslant}
\renewcommand{\le}{\leqslant}
\title[Central non-vanishing of Dirichlet $L$-functions]{Central non-vanishing of Dirichlet $L$-functions}
\author{Benjamin Durkan, Andrew Pearce-Crump}
\date{Source: arXiv:2609.06167v2 (CC BY 4.0)}
\begin{document}
\maketitle

\noindent\emph{Source and licence.} Central non-vanishing of Dirichlet $L$-functions. Version arXiv:2609.06167v2 (CC BY 4.0), distributed under the Creative Commons Attribution 4.0 International licence, as recorded in the arXiv licence metadata for this exact version and shown on the abstract page https://arxiv.org/abs/2609.06167v2. Full rights notice: licenses/arxiv-2609.06167-CC-BY-4.0.txt.\par\smallskip

\noindent\textit{Editorial scope.} Counted unit: the Lem whose source label is \texttt{lem:two-term-mixed-block-insertion} in the article (source0.tex lines 886--905, source label \texttt{lem:two-term-mixed-block-insertion}), delivered together with the article's own complete original proof of it (source0.tex lines 907--936, one \texttt{proof} environment of 30 lines). No mathematics is written, repaired or re-derived: statement and proof are the article's own text line for line. Required context retained uncounted: prop:fixed-parity-qw-setup (lines 675--711); prop:qw-preinsertion (lines 781--807). Every definition, display and label of those lines is retained unchanged. Omitted from this package and not redistributed: 1--674, 712--780, 808--885, 906--906, 937--2692. Nothing inside a retained excerpt is omitted or altered: every retained body is delivered exactly as the article's lines are written, so no omission receipt of any kind accompanies this package. Citations: the retained excerpts carry no citation command at all, so no bibliography entry of the article is transcribed and no reference is redistributed. Reference adaptation: none was needed and none was made; every reference inside the delivered unit resolves inside it, so no unresolved reference or citation is produced. Printed slips kept literal: no printed word, symbol, hypothesis or number of the counted statement or of its proof was altered. Layout and class: the article's own class is \texttt{amsart} with packages including amssymb, amsmath, amsthm, mathrsfs, enumitem, graphicx, tikz, hyperref, url; the wrapper is the lane's pinned \texttt{amsart} editorial wrapper at 10pt on A4, which restates the theorem-like environments the retained text needs and adds the article's own macro definitions that the retained text needs (\texttt{\textbackslash e}, \texttt{\textbackslash ge}, \texttt{\textbackslash le}), with extra wrapper packages (bm, bbm, dsfont, xspace, cancel, mathtools). The delivered numbering is the unit's own. Assets: no retained span contains a figure environment, a graphics-inclusion command, a PDF-page inclusion, a table or a TikZ picture; no image file is copied into this package, the package declares no asset dependency and no image file is redistributed with it. Licence: version arXiv:2609.06167v2 of this article is distributed under the Creative Commons Attribution 4.0 International licence, as recorded in the arXiv licence metadata for this exact version and shown on the abstract page https://arxiv.org/abs/2609.06167v2. A full rights notice is delivered at licenses/arxiv-2609.06167-CC-BY-4.0.txt. Characters: every retained excerpt is pure ASCII, so the delivered engine, XeLaTeX with its default Latin Modern text font, reports no missing character.
\begin{prop}\label{prop:fixed-parity-qw-setup}
Fix $\sigma>0$, lengths $0<\theta_1,\theta_2<1/2$, and real numbers $c_1,c_2$. Choose $\varepsilon_1>0$ sufficiently small in terms of $\sigma$. Let $q>1$ satisfy $q\not\equiv2\bmod4$, and let $j\in\{0,1\}$. Let the dyadic weights be supported in a fixed compact subinterval of $(0,\infty)$ and have sufficiently many uniformly bounded derivatives. Assume
$$
M_1\le q^{\theta_1+o(1)},\qquad
M_2\le q^{\theta_2+o(1)},\qquad
N_1N_2\le q^{1+o(1)},
$$
where the three $o(1)$-terms tend to zero uniformly and are bounded in absolute value by a sufficiently small constant depending on $\theta_1,\theta_2,\sigma$. Then
$$
\mathcal M_1^{(j)}(q)=c_1+c_2+O(q^{-\eta})
$$
for some $\eta>0$, depending only on $\sigma,\theta_1,\theta_2$ and the prescribed support and derivative bounds of the weights, but not on $q$; the implied constant may also depend on $c_1,c_2$.

Moreover,
$$
\begin{aligned}
\mathcal M_2^{(j)}(q)
&=
c_1^2\left(1+\frac{1}{\theta_1}\right)
+c_2^2\left(1+\frac{1}{\theta_2}\right)
+2c_1c_2+\mathcal R^{(j)}(q)
+O_{\theta_1,\theta_2}\left(\frac{\log\log q}{\log q}\right).
\end{aligned}
$$
Here $2c_1c_2$ is the contribution from $m_1m_2n_1=1$ in the two cross terms. Up to $O(q^{-\eta})$, $\mathcal R^{(j)}(q)$ is the sum over both signs, all coprime factorisations $q=q_1q_2$ with $\mu^2(q_1)=1$, and the finite dyadic subdivision. Each summand is a coefficient $O_{c_1,c_2}(1)$ times
$$
\begin{aligned}
&\frac{\mu^2(q_1)\varphi(q_2)}
{\varphi^*(q)q^{1/2}(M_1M_2N_1N_2)^{1/2}}
\sum_{\substack{m_1,m_2,n_1,n_2\\ (m_1m_2n_1n_2,q)=1\\ m_1m_2n_1>1}}
\alpha_{m_1}\beta_{m_2}
W_1\left(\frac{n_1}{N_1}\right)W_2\left(\frac{n_2}{N_2}\right)V_j(n_1n_2/q)\\
&\qquad\times\e\left(\pm \frac{n_2\overline{n_1m_1m_2q_1}}{q_2}\right).
\end{aligned}
$$
The coefficients are divisor-bounded, supported on $m_1\asymp M_1$ and $m_2\asymp M_2$, and the variables are coprime to $q$.
\end{prop}

\begin{prop}\label{prop:qw-preinsertion}
Fix $\eta_1,\eta_2>0$. Let $M_1,M_2,N_1,N_2$ be one off-diagonal dyadic block from Proposition \ref{prop:fixed-parity-qw-setup}, and assume
$$
M_1\le q^{\theta_1+o(1)},\qquad
M_2\le q^{\theta_2+o(1)},\qquad
N_1N_2\ll q^{1+\varepsilon_1},\qquad
M_1M_2\le q^{1-\eta_1}.
$$
Set $X=M_1N_2/N_1$. For a fixed factorisation $q=q_1q_2$ with $\mu^2(q_1)=1$, a fixed sign, $q_2>1$, and
$$
X\le q^{-\eta_2}
$$
with fixed $\eta_2>0$, the $k=0$ term contributes
$$
\ll_{\varepsilon_1} q^{\varepsilon_1}\frac{(M_1M_2)^{1/2}}{q}.
$$
The terms with $k\ne0$ contribute
$$
\ll_{\varepsilon_1}
q^{\varepsilon_1} q^{-3/2}M_2^{-1/2}X^{-1/2}
(Xq_2(1+X))^{3/4}A_{q_2}(2M_2)^{1/4}.
$$
If $q_2=1$, the sum over all $k\in\mathbb Z$ contributes
$$
\ll_{\varepsilon_1} q^{\varepsilon_1}\frac{(M_1M_2)^{1/2}}{q}.
$$
\end{prop}

\begin{lem}\label{lem:two-term-mixed-block-insertion}
Let $\omega_4,\omega_2\ge0$. Fix $\eta_1,\eta_2,\sigma>0$ and assume that the dyadic block lies in the complementary range
$$
\frac{N_2}{N_1}<q^\sigma\frac{M_1M_2}{q},\qquad
M_1M_2\le q^{1-\eta_1},\qquad
X=\frac{M_1N_2}{N_1}\le q^{-\eta_2}.
$$
Suppose that, for every $\varepsilon_1>0$,
$$
A_{q_2}(2M_2)\ll_{\varepsilon_1} q^{\varepsilon_1}
(q^{\omega_4}M_2^4q_2^{5/2}+q^{\omega_2}M_2^2q_2^3).
$$
If
$$
4\theta_1+6\theta_2+2\sigma+2\omega_4<3,
\qquad
2\theta_1+\theta_2+\sigma+\omega_2<1,
$$
\looseness=-1 then the fixed block $\mathcal B_\pm(q_1,q_2)$ is $O(q^{-\eta})$ for some $\eta>0$, uniformly in $q_2$, including $q_2=1$.
\end{lem}

\begin{proof}
The $k=0$ and $q_2=1$ terms are
$$
\ll q^{\varepsilon_1}\frac{(M_1M_2)^{1/2}}q\le q^{-1/2-\eta_1/2+\varepsilon_1}.
$$
 For $q_2>1$, insertion in Proposition \ref{prop:qw-preinsertion} gives for the quartic and quadratic terms, respectively,
$$
\begin{aligned}
&q^{\varepsilon_1}q^{-3/2+\omega_4/4}M_2^{1/2}X^{1/4}(1+X)^{3/4}q_2^{11/8},\\
&q^{\varepsilon_1}q^{-3/2+\omega_2/4}X^{1/4}(1+X)^{3/4}q_2^{3/2}.
\end{aligned}
$$
Since $q_2\le q$ and $X\le1$, these are at most
$$
q^{\varepsilon_1-1/8+\omega_4/4}M_2^{1/2}X^{1/4}\qquad\text{and}\qquad q^{\varepsilon_1+\omega_2/4}X^{1/4},
$$
and $X<q^{-1+\sigma}M_1^2M_2$ gives $X^{1/4}\le q^{-1/4+\sigma/4}M_1^{1/2}M_2^{1/4}$. The two contributions are therefore at most
$$
q^{\varepsilon_1-3/8+(\omega_4+\sigma)/4}M_1^{1/2}M_2^{3/4}
\qquad\text{and}\qquad
q^{\varepsilon_1-1/4+(\omega_2+\sigma)/4}M_1^{1/2}M_2^{1/4}.
$$
The resulting exponents at $M_i=q^{\theta_i+o(1)}$ are
$$
-\frac{3-(4\theta_1+6\theta_2+2\sigma+2\omega_4)}{8}+o(1),
\qquad
-\frac{1-(2\theta_1+\theta_2+\sigma+\omega_2)}{4}+o(1).
$$
Both are negative. Choose the small exponent and dyadic losses below half the smaller saving.
\end{proof}

\end{document}
